%% ============================================================
%%  appendix/A_interview_questions.tex
%%  Appendix A: Technical Interview Question Bank
%%  Credits: Ravi, PrepFusion Community
%% ============================================================

\chapter{Technical Interview Question Bank}
\label{app:questions}

\begin{protip}[How to Use This Appendix]
  Questions are organized by domain. Practice solving them on paper---many IIT
  and IISc interviews require whiteboard solutions. For digital questions, also
  draw timing diagrams and logic circuits where applicable.
\end{protip}

%% ─────────────────────────────────────────────────────────
\section{Digital Circuits}
\label{sec:iq-digital}

\subsection{Number Systems \& Representation}

\begin{enumerate}[label=\textbf{D\arabic*.}, leftmargin=*, itemsep=6pt]

  \item Convert the decimal number $156.75$ into its binary, octal, and
        hexadecimal equivalents.

  \item Convert the following decimal numbers into signed binary, octal, and
        hexadecimal using the minimum possible number of bits:
        \begin{enumerate}[label=(\alph*)]
          \item $+17$
          \item $-17$
        \end{enumerate}

  \item Explain the difference between 1's complement and 2's complement
        representations of a negative number. Why is 2's complement more widely
        used in computer systems?

  \item Explain the 2's complement representation of negative numbers. Convert
        $13$ and $-13$ into:
        \begin{enumerate}[label=(\alph*)]
          \item 8-bit signed magnitude
          \item 8-bit 1's complement
          \item 8-bit 2's complement
        \end{enumerate}

  \item What is the difference between binary and Gray code? Why is Gray code
        used in applications like rotary encoders?

\end{enumerate}

\subsection{Combinational \& Sequential Logic}

\begin{enumerate}[label=\textbf{D\arabic*.}, leftmargin=*, itemsep=6pt, start=6]

  \item Define a Flip-Flop. What is the difference between a latch and a
        flip-flop?

  \item Design a 4-bit binary multiplier circuit using basic gates. Draw the
        schematic on the whiteboard.

  \item Explain how a Wallace Tree multiplier works. How does it improve speed
        over a simple array multiplier?

  \item Sketch the timing diagram for a 2-bit multiplier with inputs $A = 11_2$
        and $B = 10_2$. Show intermediate and final outputs.

\end{enumerate}

%% ─────────────────────────────────────────────────────────
\section{Analog Circuits}
\label{sec:iq-analog}

\subsection{Oscillators \& Feedback}

\begin{enumerate}[label=\textbf{A\arabic*.}, leftmargin=*, itemsep=6pt]

  \item Draw the circuit of a Wien Bridge oscillator and explain how it sustains
        oscillations. What sets its oscillation frequency?

  \item Sketch the output waveform of an oscillator you designed. How would you
        stabilize its amplitude?

  \item Explain the role of negative feedback in an amplifier circuit. How does
        it affect gain and stability?

  \item Design an op-amp circuit with negative feedback on the whiteboard. Derive
        its closed-loop gain.

  \item What happens to an oscillator if positive feedback becomes too strong?
        Illustrate with a block diagram.

\end{enumerate}

\subsection{Control Systems (Analog Context)}

\begin{enumerate}[label=\textbf{A\arabic*.}, leftmargin=*, itemsep=6pt, start=6]

  \item What is the dominant pole concept in control systems? How does it
        influence system response?

  \item Sketch a second-order system's pole-zero plot and explain its step
        response.

  \item How would you determine the stability of a feedback system using pole
        locations? Show your approach.

\end{enumerate}

\subsection{Op-Amp Circuits}

\begin{enumerate}[label=\textbf{A\arabic*.}, leftmargin=*, itemsep=6pt, start=9]

  \item Given an op-amp circuit [provided by interviewer]: solve for the output
        voltage on the whiteboard and explain each step.

  \item Modify a given op-amp circuit to act as an integrator. Draw the new
        schematic and its output waveform for a square-wave input.

\end{enumerate}

\subsection{Differential Amplifiers}

\begin{enumerate}[label=\textbf{A\arabic*.}, leftmargin=*, itemsep=6pt, start=11]

  \item Draw a differential amplifier using MOSFETs. Calculate its differential
        gain and show the small-signal model.

  \item What happens to the CMRR if there is a mismatch in the differential pair?
        Explain with a diagram.

  \item Sketch the input and output waveforms of a differential amplifier for a
        sinusoidal differential input.

\end{enumerate}

%% ─────────────────────────────────────────────────────────
\section{Power Systems (EE Domain)}
\label{sec:iq-power}

\begin{enumerate}[label=\textbf{P\arabic*.}, leftmargin=*, itemsep=6pt]

  \item Explain the working principle of a \textbf{synchronous machine} and
        describe the role of the excitation system.

  \item What is the difference between an \textbf{induction motor} and a
        synchronous motor? Describe starting methods for induction motors.

  \item Explain \textbf{power system protection} principles: What is a
        differential relay and when is it used?

  \item Define \textbf{transient stability} in a power system. Sketch the
        equal area criterion for a fault.

  \item What are the nameplates and typical ratings of a 100~kVA,
        11~kV/415~V distribution transformer?

\end{enumerate}

%% ─────────────────────────────────────────────────────────
\section{IISc Analog Domain Interview (Sample Session)}
\label{sec:iisc-analog-session}

\begin{experience}{Anonymous Candidate}{IISc Bangalore -- Analog Domain}

  \textbf{Interviewer:} Prof.\ Aniruddhan
  \hfill\textbf{Duration:} 30--35 minutes

  \vspace{4pt}
  \textbf{Focus areas:} Oscillator design, Feedback concepts, Control Systems,
  Op-Amp circuits, Differential amplifiers.

  \vspace{4pt}
  \textbf{Style:} Intense, whiteboard problem-solving, follow-ups on written
  test mistakes.

  \begin{pfChecklist}
    \item Draw a Wien Bridge oscillator and explain oscillation sustainment.
    \item Design an op-amp circuit with negative feedback; derive closed-loop
          gain on the whiteboard.
    \item Explain the dominant pole concept and sketch a second-order system's
          step response.
    \item Draw a MOSFET differential amplifier; calculate differential gain and
          show the small-signal model.
    \item What happens to CMRR with mismatch in the differential pair?
  \end{pfChecklist}

  \textbf{Tip:} Interviewers at IISc often revisit questions you answered
  incorrectly in the written test. Prepare to solve similar problems on the spot.
\end{experience}

%% ─────────────────────────────────────────────────────────
\section{IISc Digital Domain Interview (Sample Session)}
\label{sec:iisc-digital-session}

\begin{experience}{Anonymous Candidate}{IISc Bangalore -- Digital Domain}

  \textbf{Interviewers:} Prof.\ Janakiraman \& Prof.\ Vinita Vasudevan
  \hfill\textbf{Duration:} $\approx$10 minutes

  \vspace{4pt}
  \textbf{Focus:} Multiplier circuits, domain interest.

  \begin{pfChecklist}
    \item Design a 4-bit binary multiplier using basic gates.
    \item Explain Wallace Tree multiplier and its speed advantage.
    \item Draw the timing diagram for a 2-bit multiplier: $A=11_2$, $B=10_2$.
  \end{pfChecklist}

  \textbf{Note:} The panel also asked about the candidate's genuine interest in
  the digital domain vs.\ analog, and requested an example of a digital circuit
  they had worked on.
\end{experience}
